\documentclass[11pt,a4paper]{article}

\usepackage[utf8]{inputenc}
\usepackage[T1]{fontenc}
\usepackage[margin=2.5cm]{geometry}
\usepackage{graphicx}
\usepackage{booktabs}
\usepackage{amsmath}
\usepackage{float}
\usepackage{siunitx}
\usepackage{xcolor}
\usepackage[hidelinks]{hyperref}
\usepackage{caption}

\graphicspath{{../results/plots/}}

\title{\textbf{Parallel Traffic-Light Optimization on an HPC Cluster}\\
\large Large Scale Computing --- Project Report}
\author{Grzegorz Mróz and Katarzyna Kucharzyk}
\date{\today}

\begin{document}
\maketitle

\begin{abstract}
This project uses traffic-light timing optimization as a realistic, embarrassingly
parallel workload to demonstrate core High Performance Computing (HPC) concepts:
SLURM job arrays, parallel scaling, search-strategy comparison, parameter
sensitivity analysis, and an idealized Monte Carlo baseline. Every candidate
traffic-light configuration is evaluated by an independent simulation, which makes
the workload trivially parallelizable across the nodes of the PLGrid/Ares cluster.
We report four experiments. The scaling study confirms near-linear speedup
(efficiency $\geq 91\%$ up to $100$ workers). An evolutionary algorithm (EA)
outperforms random and grid search at an equal simulation budget, reducing the
mean objective by roughly $6\times$ and cutting average vehicle waiting time by
$10.7\%$ relative to a hand-crafted fixed-time baseline. A sensitivity sweep shows that, for a fixed
budget, more generations beat larger populations. Finally, a Monte Carlo
$\pi$-estimator on $5\times10^{7}$ samples illustrates perfect parallelism with
zero inter-task communication.
\end{abstract}

\section{Introduction}

Optimizing the timing of traffic lights is computationally expensive because each
candidate configuration must be evaluated in a traffic simulation. Since the
configurations are mutually independent, the search is an \emph{embarrassingly
parallel} problem: thousands of simulations can run simultaneously with no
communication between them. This makes traffic-light optimization an ideal
vehicle for studying HPC batch processing.

The goal of this project is \textbf{not} to build a production traffic controller,
but to use a scalable prototype to demonstrate and measure HPC properties:
how the workload scales with the number of workers, how different optimization
strategies trade off quality against compute budget, and how an idealized parallel
workload behaves as a reference point.

\section{Methodology}

\subsection{Traffic-light model and objective}

Each intersection is controlled by a fixed-time signal program described by a JSON
configuration with the following parameters: \texttt{cycle\_length}, North--South
green time (\texttt{green\_ns}), East--West green time (\texttt{green\_ew}),
\texttt{yellow\_time}, \texttt{all\_red\_time}, and a coordination \texttt{offset}.
All configurations satisfy the timing constraint
\[
\texttt{green\_ns} + \texttt{green\_ew} + 2\cdot\texttt{yellow\_time}
+ 2\cdot\texttt{all\_red\_time} \;\leq\; \texttt{cycle\_length}.
\]
The optimization objective is to \textbf{minimize the average vehicle waiting time}
($\texttt{avg\_waiting\_time}$, in seconds) over a one-hour simulated horizon on a
$3\times3$ grid of intersections under medium traffic demand.

\subsection{Simulation backends}

The architecture follows one strict rule: \emph{nothing outside the simulator
module knows how a backend works internally}. Every backend consumes one JSON
config and emits exactly one CSV row with a fixed metric schema. Three backends
are provided:
\begin{itemize}
  \item \texttt{python} --- a discrete-time queueing model of the intersections.
        Fast ($\sim$\SI{7}{\milli\second} per simulation) and dependency-free.
  \item \texttt{sumo} --- a microscopic simulation using the SUMO traffic
        simulator ($\sim$\SI{1.4}{\second} per simulation). Used here for
        visualization.
  \item \texttt{mc} --- a Monte Carlo $\pi$-estimator used as an idealized
        embarrassingly-parallel reference.
\end{itemize}
Because all backends share the same input/output contract, the SLURM workflow,
aggregation, and analysis code are backend-agnostic. The four headline
experiments were executed on Ares with the \texttt{python} backend (SUMO was not
available in the cluster module environment); the SUMO backend was used locally
to render the traffic animation.

\subsection{HPC execution model}

Each simulation is one task of a SLURM \emph{job array}. Concurrency is controlled
with the array throttle (\texttt{--array=0-P\%N}), which limits how many of the
$P$ tasks run simultaneously to $N$. Results are written as independent CSV files
and aggregated afterwards, so there is no shared state and no synchronization
during evaluation. The EA introduces the only synchronization point: a barrier
between generations, where the next population is bred from the evaluated parents.

\section{Experiment 1: Parallel Scaling}

\paragraph{Setup.} A fixed workload of $P=100$ configurations is evaluated while
varying the number of concurrent workers $N \in \{1,2,5,10,20,50,100\}$ (strong
scaling). Each task records its own runtime. The wall-clock time for $N$ workers
is modelled as the makespan of scheduling the measured per-task runtimes onto $N$
workers (longest-processing-time greedy schedule); the sequential baseline $T(1)$
is the sum of all task runtimes. Speedup is $S(N)=T(1)/T(N)$ and efficiency is
$E(N)=S(N)/N$.

\begin{table}[H]
\centering
\caption{Strong-scaling results ($P=100$ tasks).}
\begin{tabular}{rrrr}
\toprule
Workers $N$ & Wall time $T(N)$ [s] & Speedup $S(N)$ & Efficiency $E(N)$ \\
\midrule
1   & 0.6875 & 1.00  & 100.0\% \\
2   & 0.3438 & 2.00  & 100.0\% \\
5   & 0.1376 & 5.00  & 99.9\%  \\
10  & 0.0688 & 9.99  & 99.9\%  \\
20  & 0.0345 & 19.93 & 99.6\%  \\
50  & 0.0139 & 49.46 & 98.9\%  \\
100 & 0.0075 & 91.67 & 91.7\%  \\
\bottomrule
\end{tabular}
\end{table}

\begin{figure}[H]
\centering
\includegraphics[width=\linewidth]{scaling.png}
\caption{Strong scaling. Left: measured speedup tracks the ideal linear curve
almost exactly. Right: parallel efficiency stays above $98\%$ up to $N=50$ and
only drops to $91.7\%$ at $N=100$, where granularity (one task per worker) makes
load imbalance unavoidable.}
\end{figure}

\paragraph{Discussion.} The workload exhibits essentially ideal strong scaling, as
expected for an embarrassingly parallel problem: there is no inter-task
communication, so the only loss of efficiency comes from load granularity. At
$N=100$ each worker handles a single task, so the longest task dictates the
makespan and efficiency falls to $91.7\%$ --- a textbook illustration of why the
number of tasks should comfortably exceed the number of workers.

\section{Experiment 2: EA vs.\ Random vs.\ Grid Search}

\paragraph{Setup.} Three search strategies are compared: \emph{random search}
(300 random configs), \emph{grid search} (a regular lattice over cycle length and
green splits, 72 points --- the natural resolution of the grid before
combinatorial blow-up), and an \emph{evolutionary algorithm} (population of 50
over 6 generations, i.e.\ 300 evaluations). Random search and the EA thus operate
on an equal budget of 300 evaluations, while grid search samples the space more
sparsely. Quality is the best/mean/worst $\texttt{avg\_waiting\_time}$. As an
absolute reference we also report a hand-crafted \emph{fixed-time baseline}
(\SI{60}{\second} cycle, $30/25$ green split, zero offset), evaluated on the same
backend and demand.

\begin{table}[H]
\centering
\caption{Search-strategy comparison (avg.\ waiting time, seconds; lower is
better). EA statistics describe the final-generation population of 50
individuals; its total budget is 300 evaluations.}
\begin{tabular}{lrrrr}
\toprule
Strategy & Evaluations & Best & Mean & Worst \\
\midrule
Fixed-time baseline & 1 & 12.25 & --- & --- \\
Random search & 300 & 11.03 & 212.87 & 1070.36 \\
Grid search   & 72  & 11.55 & 138.40 & 710.20  \\
Evolutionary algorithm & 300 (50 $\times$ 6 gen.) & \textbf{10.94} & \textbf{33.94} & \textbf{366.89} \\
\bottomrule
\end{tabular}
\end{table}

\begin{figure}[H]
\centering
\includegraphics[width=\linewidth]{search_comparison.png}
\caption{Left: best/mean/worst objective per strategy. Middle: cumulative
distribution of solution quality --- the EA curve rises fastest, i.e.\ it produces
far fewer poor solutions. Right: best-so-far learning curve against the number of
simulations.}
\end{figure}

\paragraph{Discussion.} All three strategies find a similar \emph{best} solution
($\approx$\SI{11}{\second}), because good configurations are not especially rare.
The decisive difference is the \emph{distribution}: the EA's mean objective
(\SI{33.9}{\second}) is about $6\times$ lower than random search
(\SI{212.9}{\second}) and $4\times$ lower than grid search, and its worst case is
roughly $3\times$ better. The CDF makes this explicit --- the EA concentrates its
budget in the high-quality region instead of wasting evaluations on poor
candidates. For a fixed compute budget, the EA therefore extracts substantially
more value, which is the central motivation for using guided search on HPC.
Grid search received fewer evaluations than the other two methods, but this only
strengthens the conclusion for random search (equal budget, $6\times$ worse mean)
and does not change the ranking: the EA dominates on every reported statistic.

\paragraph{Practical payoff.} Measured against the fixed-time baseline
(\SI{12.25}{\second} average waiting under identical demand --- 3262 vs.\ 3268
vehicle arrivals), the best evolved configuration achieves
\SI{10.94}{\second}, i.e.\ a $\mathbf{10.7\%}$ \textbf{reduction in average
vehicle waiting time}. This is the bottom line of the optimization: a sensibly
hand-tuned signal plan is already decent, yet a few CPU-minutes of parallel
guided search buy a double-digit relative improvement.

\section{Experiment 3: EA Parameter Sensitivity}

\paragraph{Setup.} The EA budget is fixed at 300 evaluations and split between
population size and number of generations in seven ways, from $300\times1$
(pure random) to $15\times20$ (deep, narrow search).

\begin{table}[H]
\centering
\caption{EA sensitivity at a fixed budget of 300 evaluations.}
\begin{tabular}{rrr}
\toprule
Population & Generations & Best avg.\ waiting time [s] \\
\midrule
300 & 1  & 11.03 \\
150 & 2  & 11.45 \\
100 & 3  & 11.47 \\
60  & 5  & 10.88 \\
50  & 6  & 11.51 \\
30  & 10 & 10.61 \\
15  & 20 & \textbf{10.55} \\
\bottomrule
\end{tabular}
\end{table}

\begin{figure}[H]
\centering
\includegraphics[width=\linewidth]{sensitivity.png}
\caption{Left: best-fitness convergence per configuration. Middle and right: final
best objective as a function of population size and of the number of generations.}
\end{figure}

\paragraph{Discussion.} At a fixed budget there is a clear exploration vs.\
exploitation trade-off. Configurations with a single generation ($300\times1$,
equivalent to random search) and very large populations leave quality on the
table, whereas allocating the budget to \emph{more generations}
($15\times20$, $30\times10$) yields the best results. This indicates that
iterative selection pressure --- not raw population diversity --- drives
improvement for this problem. The differences are modest (all within
$\approx$\SI{1}{\second}), reflecting that the EA used here has no elitism, so the
incumbent best can be lost between generations; adding elitism is the natural next
improvement.

\section{Experiment 4: Monte Carlo $\pi$ --- An Idealized Parallel Baseline}

\paragraph{Setup.} As a reference for ``perfect'' parallelism, 500 independent
tasks each draw $10^{5}$ random points in the unit square and count those inside
the quarter circle, giving $\pi \approx 4\,(\text{hits}/\text{samples})$. The
combined estimate uses $5\times10^{7}$ samples. Unlike the EA, there is no
synchronization barrier whatsoever.

\begin{figure}[H]
\centering
\includegraphics[width=\linewidth]{mc_convergence.png}
\caption{Left: running estimate of $\pi$ converging to the true value within its
theoretical $95\%$ confidence band. Middle: absolute error vs.\ number of tasks on
log--log axes, following the expected $1/\sqrt{N}$ trend. Right: speedup ---
Monte Carlo is linear (no communication) while the EA's per-generation barrier
limits its scaling.}
\end{figure}

\paragraph{Results.} The aggregate estimate is $\pi \approx 3.141626$ versus the
true $3.141593$, an absolute error of $3.3\times10^{-5}$, consistent with the
$1/\sqrt{N}$ law. This experiment isolates the ideal case: speedup is perfectly
linear because tasks never communicate, in contrast to the EA, whose generation
barrier introduces a (small) serial fraction and caps achievable speedup ---
a concrete illustration of Amdahl's law.

\section{Visualization}

The SUMO backend produces per-vehicle traces that are rendered into an animation
of the optimized $3\times3$ grid (\texttt{results/plots/anim\_best.gif}).
Figure~\ref{fig:anim} shows a representative frame: cyan markers indicate active
green phases and red markers show queued vehicles waiting at red lights.

\begin{figure}[H]
\centering
\includegraphics[width=0.65\linewidth]{anim_best_frame.png}
\caption{Snapshot of the SUMO microsimulation for the best configuration
($t=\SI{900}{\second}$, 58 vehicles in the network).}
\label{fig:anim}
\end{figure}

\section{Conclusions}

\begin{itemize}
  \item \textbf{Embarrassingly parallel, near-ideal scaling.} The simulation
        workload scales almost linearly, with efficiency above $98\%$ up to 50
        workers; the only loss comes from task granularity at $N=100$.
  \item \textbf{Guided search wins at equal cost.} The EA delivers a $\sim6\times$
        lower mean objective and a markedly better worst case than random or grid
        search for the same simulation budget.
  \item \textbf{Optimization pays off in traffic terms.} The best evolved signal
        plan reduces average vehicle waiting time by $10.7\%$ relative to a
        hand-crafted fixed-time baseline ($12.25 \rightarrow$ \SI{10.94}{\second}).
  \item \textbf{Spend budget on iterations.} Under a fixed budget, more
        generations outperform larger populations, highlighting the value of
        selection pressure.
  \item \textbf{Synchronization has a cost.} The communication-free Monte Carlo
        baseline scales perfectly, whereas the EA's generation barrier introduces
        a serial fraction --- a practical demonstration of Amdahl's law.
\end{itemize}

Overall, the project shows how a realistic optimization problem maps cleanly onto
HPC batch infrastructure and how the choice of algorithm, not just raw compute,
determines the value obtained from a fixed allocation of CPU hours.

\end{document}
